\section*{Capítulo \ref{ch:b2:weather-systems} --- Sistemas meteorológicos}

\begin{solution}{exo:b2:weather-systems:1}
(a) $f = 2 \times 7.292 \times 10^{-5} \times \sin\ang{50} = \qty{1.117e-4}{s^{-1}}$. (b) Con
$\bar T = \qty{278}{K}$, $\tan\alpha = 1.117 \times 10^{-4} \times 278 \times 15/(9.81 \times 6)
= 0.00791$, aproximadamente 1 por cada 126. (c) $1/0.00791 = 126$: a unos \qty{126}{km} por detrás del frente. (d) La
pendiente es inversamente proporcional al contraste de temperatura: se reduce a la mitad, a 1 por cada 253.
\end{solution}

\begin{solution}{exo:b2:weather-systems:2}
(a) $6 \times 150 = \qty{900}{km}$ por delante. (b) $900/25 = \qty{36}{h}$, un día y medio. (c) Primero
los cirros altos, luego los cirrostratos, luego los altostratos; la nubosidad se espesa y baja, y la lluvia
empieza a caer de la nube baja y espesa cuando el frente se acerca.
\end{solution}

\begin{solution}{exo:b2:weather-systems:3}
Los símbolos apuntan hacia donde se mueve el frente, y hacia el aire más cálido por delante de un frente frío o el
aire más frío por delante de un frente cálido. (a) Un frente frío: el aire frío está al noroeste, por detrás de
él; se mueve hacia el sureste. (b) Un frente cálido: el aire frío está al norte, por delante de él; se mueve hacia el
norte. (c) Un \omterm{def:b2:weather-systems:fronts}{frente ocluido}, que une la baja con el punto triple donde empiezan los frentes frío y
cálido; hay aire frío a ambos lados, con el aire cálido levantado por encima; se mueve hacia donde apuntan sus
símbolos. (d) Un \omterm{def:b2:weather-systems:fronts}{frente estacionario}: los triángulos apuntan al sur, hacia el aire cálido, así que el aire frío
está al norte; apenas se mueve.
\end{solution}

\begin{solution}{exo:b2:weather-systems:4}
(a) $\beta = 2 \times 7.292 \times 10^{-5} \times \cos\ang{45}/(6.371 \times 10^6) =
\qty{1.62e-11}{m^{-1}.s^{-1}}$. (b) Para \qty{4000}{km}, $k = \qty{1.571e-6}{m^{-1}}$, $\beta/k^2 =
\qty{6.6}{m/s}$ y $c = 12 - 6.6 = \qty{5.4}{m/s}$; para \qty{8000}{km}, $\beta/k^2 =
\qty{26.3}{m/s}$ y $c = \qty{-14.3}{m/s}$. (c) La onda corta se desplaza hacia el este, \qty{470}{km} al día;
la larga se desplaza hacia el oeste, \qty{1240}{km} al día, contra el viento.
\end{solution}

\begin{solution}{exo:b2:weather-systems:5}
(a) $c = 0$ cuando $U = \beta/k^2$, así que $k = \sqrt{\beta/U}$ y $\lambda = 2\pi/k =
2\pi\sqrt{U/\beta}$. (b) $2\pi\sqrt{12/(1.62 \times 10^{-11})} = \qty{5410}{km}$. (c) El círculo
de latitud \ang{45} mide $2\pi \times 6371 \times \cos\ang{45} = \qty{28300}{km}$: unas 5.2
ondas, así que cinco en la práctica. (d) $\lambda$ crece como $\sqrt U$: un viento del oeste más fuerte alarga las
ondas estacionarias, y caben menos alrededor del hemisferio.
\end{solution}

\begin{solution}{exo:b2:weather-systems:6}
(a) En $C_+ = (\sqrt{72}, \sqrt{72}, 27)$: $\dot x = 10(\sqrt{72} - \sqrt{72}) = 0$, $\dot y =
\sqrt{72}(28 - 27) - \sqrt{72} = 0$, $\dot z = 72 - \tfrac83 \times 27 = 72 - 72 = 0$.
(b) $\mathrm{e}^{-13.67} = 1.16 \times 10^{-6}$: el volumen se contrae casi un millón de veces en una
unidad de tiempo. (c) Un volumen puede contraerse mientras se estira en una dirección y se comprime con mucha más
fuerza en las otras: la distancia según la dirección de estiramiento crece, y dos trayectorias próximas
se separan a lo largo de ella, mientras el volumen entre ellas se anula.
\end{solution}

\begin{solution}{exo:b2:weather-systems:7}
(a) $\ln 2/0.906 = 0.765$ unidades de tiempo. (b) $\ln(10/10^{-8})/0.906 = \ln(10^9)/0.906 = 22.9$ unidades
de tiempo. (c) $\ln 1000/0.906 = 7.6$ unidades de tiempo: mejorar mil veces el estado inicial
solo compra un tercio del tiempo hallado en (b), ya que $\ln 10^3/\ln 10^9 = 1/3$.
\end{solution}

\begin{solution}{exo:b2:weather-systems:8}
(a) $(303 - 195)/195 = 0.554$ y $V_{\max} = \sqrt{0.9 \times 0.554 \times 15\,000} =
\qty{86.5}{m/s}$. (b) Muy por encima de los \qty{70}{m/s} de la categoría 5, cerca de los \qty{86}{m/s} que
han superado cinco tormentas recientes. (c) Las tormentas rara vez permanecen mucho tiempo en el estado estacionario que supone
la teoría; la cizalladura vertical del viento, el aire seco que entra en la tormenta y el aire que desciende a su alrededor, que
la teoría deja fuera, las debilitan.
\end{solution}

\begin{solution}{exo:b2:weather-systems:9}
(a) $(303 - 195)/303 = 0.356$. (b) $(303 - 195)/195 = 0.554$; el cociente es $0.554/0.356 =
303/195 = 1.554$. (c) En una máquina de Carnot el trabajo sale de la máquina. En un ciclón tropical el
trabajo se disipa por rozamiento cerca de la superficie del mar, y ese calor calienta el aire que entra a la
temperatura $T_s$, donde la máquina toma su calor: se vuelve a usar, y la máquina se comporta
como si fuera más eficiente en el factor $T_s/T_o$.
\end{solution}

\begin{solution}{exo:b2:weather-systems:10}
(a) La deducción de la \cref{prop:b2:weather-systems:slope} sin aproximación da
$\mathrm{d}z/\mathrm{d}x = f(\rho_c v_c - \rho_w v_w)/[g(\rho_c - \rho_w)]$. Un fluido más denso
que descansa sobre otro más ligero es inestable: la menor perturbación lo hace volcar. En un frente en
reposo, el aire frío y más denso debe estar por tanto debajo, que es la cuña de la definición; entonces
$\rho_c > \rho_w$, la superficie se eleva hacia el aire frío y $\rho_c v_c > \rho_w v_w$.
(b) $\zeta = \partial v/\partial x = 10/(50 \times 10^3) = \qty{2e-4}{s^{-1}}$, unas $1.9f$ con
$f = \qty{1.03e-4}{s^{-1}}$. (c) Positiva en el hemisferio norte: ciclónica, y tan grande como
el propio parámetro de Coriolis; los frentes son bandas de fuerte vorticidad ciclónica.
\end{solution}

\begin{solution}{exo:b2:weather-systems:11}
(a) $\log_2(6/0.5) = \log_2 12 = 3.58$ duplicaciones, $3.58 \times 2.4 = 8.6$~días. (b) Un
día más de predicción exige un tiempo de duplicación menos de crecimiento antes de que el error sea inútil: el
error inicial debe reducirse en un factor $2^{1/2.4} = 1.33$ cada década. (c) Una reducción a la décima parte equivale a
$\log_2 10 = 3.32$ duplicaciones; si solo gana de 4 a 5 días, los errores pequeños se duplican en $4/3.32 = 1.2$
a $5/3.32 = 1.5$ días: los errores pequeños crecen más deprisa que los mayores medidos en las
predicciones de la década de 1980, por el calor latente liberado en la convección a pequeña escala.
\end{solution}

\begin{solution}{exo:b2:weather-systems:12}
(a) A \qty{30}{\celsius}, $e_s = \qty{43.5}{hPa}$ y $q_s^* = \qty{27.1}{g/kg}$; el aire a
\qty{29}{\celsius} tiene $e = 0.8 \times 41.0 = \qty{32.8}{hPa}$ y $q_a = \qty{20.4}{g/kg}$;
$h_o^* - h^* = 1004.5 \times 1 + 2.501 \times 10^6 \times 0.0067 = \qty{17.8}{kJ/kg}$, frente a
\qty{16.8}{kJ/kg}. (b) $(303.15 - 200)/200 = 0.516$ y $V_{\max} = \sqrt{0.9 \times 0.516
\times 17\,800} = \qty{90.9}{m/s}$, frente a \qty{87.9}{m/s}. (c) Alrededor de un \qty{3.4}{\percent} por
kelvin. El desequilibrio aumenta un \qty{6}{\percent}, lo que añade un \qty{3}{\percent} a la velocidad;
el factor $(T_s - T_o)/T_o$ aumenta un \qty{1}{\percent}, lo que añade un \qty{0.5}{\percent}: el
combustible, la humedad que puede aportar el mar más cálido, es lo que más contribuye, por el aumento exponencial de
$e_s$.
\end{solution}

\begin{solution}{pb:b2:weather-systems:1}
\textbf{1.} $(480 - 390)/6 = \qty{15}{km/h}$.
\textbf{2.} $390/15 = \qty{26}{h}$ después de las 12 UTC: a las 14 UTC del día siguiente.
\textbf{3.} $f = \qty{1.117e-4}{s^{-1}}$.
\textbf{4.} Con $\bar T = \qty{281}{K}$, $\tan\alpha = 1.117 \times 10^{-4} \times 281 \times
16/(9.81 \times 8) = 0.0064$, aproximadamente 1 por cada 156.
\textbf{5.} El frente está entonces a $6 \times 15 = \qty{90}{km}$ más allá de la ciudad; la superficie está
a $90\,000 \times 0.0064 = \qty{576}{m}$ por encima de ella.
\textbf{6.} $1500/0.0064 = \qty{234}{km}$ por detrás del frente, alcanzados $234/15 = \qty{15.6}{h}$ después
del paso.
\textbf{7.} Cúmulos y cumulonimbos, chubascos, quizá tormenta, rachas y una caída repentina de la
temperatura; la presión, que baja mientras el frente se acerca, empieza a subir a su paso; después
un aire más fresco y más limpio.
\textbf{8.} $\beta = 2 \times 7.292 \times 10^{-5} \times \cos\ang{50}/(6.371 \times 10^6) =
\qty{1.47e-11}{m^{-1}.s^{-1}}$.
\textbf{9.} $k = 2\pi/(5 \times 10^6) = \qty{1.257e-6}{m^{-1}}$ y $\beta/k^2 = \qty{9.3}{m/s}$.
\textbf{10.} $c = 20 - 9.3 = \qty{10.7}{m/s}$, hacia el este.
\textbf{11.} $10.7 \times 86\,400 = \qty{924}{km}$ hacia el este.
\textbf{12.} $\lambda_s = 2\pi\sqrt{20/(1.47 \times 10^{-11})} = \qty{7330}{km}$; el círculo de
latitud \ang{50} mide $2\pi \times 6371 \times \cos\ang{50} = \qty{25700}{km}$, unas 3.5
ondas así.
\textbf{13.} Por delante de ella, en su lado este, donde las ondas cortas inducen aire ascendente y la
presión en superficie baja.
\textbf{14.} $\lambda = \ln(3.9 \times 10^{-3}/5 \times 10^{-7})/(14 - 4) = \ln 7800/10 = 0.896$.
\textbf{15.} $\ln 2/0.896 = 0.77$ unidades de tiempo.
\textbf{16.} Un \qty{1}{\percent} por debajo de 0.906: la media de 200 pares, un tiempo finito y lecturas hechas en
una gráfica.
\textbf{17.} Una vez que las dos pasadas están muy separadas, ocupan puntos sin relación del atractor, y
su distancia no puede superar su tamaño: el crecimiento se satura.
\textbf{18.} $\ln 10/0.896 = 2.6$ unidades de tiempo más tarde.
\textbf{19.} $\log_2(3/0.1) = \log_2 30 = 4.91$ duplicaciones.
\textbf{20.} $4.91 \times 2.4 = 11.8$~días.
\textbf{21.} Un tiempo de duplicación: 2.4~días.
\textbf{22.} $4.91 \times 5 = 24.5$~días: la estimación antigua era demasiado optimista.
\textbf{23.} Con un tiempo de duplicación de 2.4 días, un error diez veces menor ganaría $3.32 \times 2.4
= 8.0$~días, no de 4 a 5: los errores pequeños crecen más deprisa que una duplicación cada 2.4 días, y una
estimación basada en ese tiempo de duplicación sobrevalora el horizonte.
\textbf{24.} Esta predicción podría ser útil durante unos \textbf{12 días} como mucho, y en la práctica
menos.
\end{solution}
